% =============================================================================
% Chapter 2 --- Background
% Draws strictly from the four PDFs in report/Background_Papers/:
%   - Verma et al. 2019 (IEEE SSC Magazine)
%   - Sebastian et al. 2020 (Nature Nanotechnology)
%   - Ielmini & Wong 2018 (Nature Electronics)
%   - Wang et al. 2022 (arXiv:2202.10228, Y-Flash memristor, Technion)
% =============================================================================

\chapter{Background}
\label{ch:background}

This chapter provides the technical context needed to appreciate the role
of the controller designed in this project. It is organized around four
themes drawn directly from the cited literature: the energetic argument
for in-memory computing, the role of resistive switching devices as
programmable synaptic elements, the write--verify problem that such devices
impose on their host digital control, and the Y-Flash floating-gate
memristor that is relevant to the Technion device and circuit ecosystem
in which this project is set.

\section{In-Memory Computing and the Memory Wall}
\label{sec:bg-imc}

Contemporary computing platforms, including those specialized for machine
learning, are limited by the cost of moving data between memory and
processing units rather than by the cost of arithmetic itself. As Sebastian
et~al.~\cite{Sebastian2020IMC} observe, even at relatively mature technology
nodes the energy cost of multiplying two numbers is orders of magnitude
smaller than the energy cost of fetching those numbers from memory.
Conventional mitigations, including parallel architectures such as GPUs and
application-specific digital
accelerators~\cite{Ielmini2018RRAM, Sebastian2020IMC}, reduce the cost of
computation but leave the fundamental data-movement penalty largely intact.

The in-memory computing paradigm takes a different route: instead of moving
data to the processor, it embeds the computation inside a memory array and
reads out a result that aggregates many per-device contributions in
parallel~\cite{Sebastian2020IMC, Verma2019IMC}. For the dominant kernel of
ANN inference --- high-dimensional matrix--vector multiplication --- there
is a natural structural alignment between a 2D array of bit cells and the
dataflow of the kernel: weights are stored in the cells, inputs are applied
on one dimension of the array, and per-cell contributions combine physically
along the orthogonal dimension. Recent prototype silicon demonstrates that
this approach can yield order-of-magnitude improvements in both energy
efficiency and throughput for matrix-vector multiplication relative to
conventional digital accelerators~\cite{Verma2019IMC}.

However, fitting computation inside a constrained bit-cell array imposes
tangible system-level constraints. Analog computation in the array must be
bridged to conventional digital accelerators for programmability and
integration; the array must be virtualizable so that software can be mapped
onto it~\cite{Verma2019IMC}; and the peripheral circuitry, control logic,
and interconnect must be co-designed with the array so that the benefits
of in-memory computation are not consumed by inefficiencies at the
periphery~\cite{Sebastian2020IMC}. The controller designed in this report
is precisely the digital component that lives at this interface.

\section{Resistive Switching Devices as Synaptic Elements}
\label{sec:bg-rram}

At the device level, a wide family of emerging non-volatile memories has
been proposed as candidates for IMC storage elements: resistive switching
RAM (RRAM), phase-change memory (PCM), magnetoresistive RAM (MRAM), and
ferroelectric RAM (FeRAM), among others~\cite{Ielmini2018RRAM}. These
devices share a common property that is important for ANN acceleration:
their conductance can be programmed to a continuum of intermediate
values, not merely a binary state, and it persists without power. A
programmable conductance can therefore directly represent a synaptic
weight in an analog fashion, and a crossbar array of such devices
implements MVM in one read cycle by summing per-column currents induced by
per-row input voltages~\cite{Ielmini2018RRAM, Sebastian2020IMC}.

As Ielmini and Wong~\cite{Ielmini2018RRAM} discuss in their review of
in-memory computing with resistive switching devices, this arrangement
blurs the boundary between memory and computation, and it is supported by
a growing corpus of prototype chips and system demonstrations. The same
authors, however, also emphasize that resistive switching devices exhibit
non-idealities --- including device-to-device variability, cycle-to-cycle
variability, retention drift, and finite endurance --- that distinguish
them from conventional digital memory and that have direct implications
for how a host controller must interact with the array.

\section{The Write--Verify Problem}
\label{sec:bg-writeverify}

A fundamental consequence of the device non-idealities summarized above is
that a single unconditional \emph{write} operation cannot be relied upon
to place a cell at its intended conductance level. Sebastian
et~al.~\cite{Sebastian2020IMC} describe how programming resistive devices
typically requires an iterative strategy in which programming pulses are
interleaved with read operations that check the cell's state against the
target value. Pulse amplitude, duration, and polarity may be adapted
between iterations to move the cell toward the intended level, and the
sequence is terminated either when the readback matches the target or when
a retry budget is exhausted.

This \emph{program--verify--(re-program or erase)} loop is the primary
reason that a controller for an IMC-based ANN accelerator is substantially
more complex than a conventional write-path controller. It must manage
the pulse train for each programming attempt, schedule a read-back phase,
compare the observed weight against an expected value, and branch into
either a retry on the same cell, an erase followed by reprogramming, or
a declaration that the cell cannot be brought to the target in the
current configuration. In the present project, exactly this control
structure is implemented as a hierarchical finite-state machine with
explicit PROG, VERIFY, and ERASE sub-FSMs, and it is the central object
of \Cref{ch:rtl}.

\section{Y-Flash Floating-Gate Memristors}
\label{sec:bg-yflash}

Among emerging non-volatile technologies, the Y-Flash floating-gate
memristor~\cite{Wang2022YFlash} is of particular interest in the
Technion device and circuit ecosystem in which this project is situated.
Y-Flash devices are fabricated in a standard CMOS process with a
single-polysilicon floating gate, and they operate in a two-terminal
configuration that is analogous to other emerging memristive
technologies~\cite{Wang2022YFlash}. Compared with filamentary RRAM or
phase-change devices, Y-Flash cells can achieve a large number of
fine-tuned intermediate conductance states in the subthreshold region
with low readout currents (on the order of 1\,nA to 5\,\si{\micro\ampere}),
high retention, low cycle-to-cycle variability, and high
yield~\cite{Wang2022YFlash}, all of which are desirable properties for a
large-scale synaptic array.

Arranged in a crossbar, Y-Flash devices can accelerate the frequent and
expensive vector--matrix multiplications that dominate neuromorphic
inference~\cite{Wang2022YFlash}. The device supports explicit
\emph{program}, \emph{erase}, and \emph{readout} operations, each of which
must be issued by the host under an appropriate pulse shape. Wang
et~al.~\cite{Wang2022YFlash} provide a physical-based compact model that
describes program, erase, and readout dynamics and that is intended to be
co-designed with the peripheral CMOS control circuitry --- reinforcing
the more general point made by Sebastian et~al.~\cite{Sebastian2020IMC}
that the value delivered by an IMC array is determined jointly by the
device, the array, and the control logic.

\section{Implications for the Present Project}
\label{sec:bg-implications}

The background reviewed in this chapter motivates three architectural
choices that are visible in the RTL described in the remainder of the
report:

\begin{itemize}
    \item A \emph{packed, one-hot} address and data representation that
          cleanly maps to a tiled crossbar-like core structure, so that
          a programming or read command uniquely selects one cell per
          block and sub-block.
    \item A \emph{pulse-train} command interface on three bits
          (\signame{pulses}[2:0]) that encodes read, program, erase, and
          inference modes as distinct time-multiplexed bursts, rather
          than as arbitrary bit patterns; this corresponds to the
          device-level distinction between the different operations
          discussed in~\cite{Wang2022YFlash, Sebastian2020IMC}.
    \item A \emph{closed-loop} control policy with explicit PROG, VERIFY,
          and ERASE sub-FSMs and a bounded retry budget, reflecting the
          write--verify requirement imposed by non-ideal resistive
          devices~\cite{Ielmini2018RRAM, Sebastian2020IMC}.
\end{itemize}

With this background in place, \Cref{ch:architecture} begins the
description of the actual RTL by presenting the three-module stack and
the host-to-core word format that binds them together.
